\section{Introduction}
  \label{sec:introduction}

  Gravity is traditionally described as a geometric phenomenon.
  In general relativity, the gravitational interaction is encoded in the curvature of spacetime, and the Newtonian
  potential arises as a weak-field limit of the Einstein equations.
  This geometric formulation is remarkably successful.
  It nevertheless leaves open the conceptual question of whether gravitational dynamics must be fundamental, or
  whether it may arise as an effective response of a more primitive operator-based structure~\cite{Cosmochrony}.

  In parallel, operator-theoretic and spectral approaches have clarified how elliptic operators encode geometric and
  physical information.
  Functional determinants, resolvents, and heat-kernel asymptotics play a central role in quantum field theory,
  statistical mechanics, and spectral geometry~\cite{Vassilevich2003,ParkerToms2009}.
  Variations of logarithmic determinants generate non-local responses through the inverse operator.
  This suggests a bridge between entropy-like functionals and long-range interaction.

  In a previous analysis~\cite{Beau2026a}, we investigated a minimal mechanism by which a Newtonian interaction
  can emerge from the
  variation of a projective entropy functional in three spatial dimensions.
  For a positive Laplace-type operator \(A\), the functional
  \begin{equation}
    S_\Pi[A] = \frac12 \log \det\nolimits' A
  \end{equation}
  has the weak-perturbation variation
  \begin{equation}
    \delta S_\Pi = \frac12 \Tr\!\left(A^{-1}\delta A\right),
  \end{equation}
  so that the Green kernel governs the spatial structure of the response.
  In three dimensions, the Green kernel of a Laplace-type operator exhibits a universal \(1/r\) decay.
  Consequently, the entropy response inherits a Newtonian long-range profile without assuming a force law.

  The goal of the present paper is to develop a covariant extension of this mechanism and to identify, in a controlled
  infrared regime, the leading local geometric content of the \emph{renormalized} metric variation.
  We do not claim an exact equality between the metric variation of the bare determinant and Einstein's equations.
  We also do not discard non-local contributions, higher-derivative curvature terms, or the conformal anomaly.
  Instead, our central claim is the following.

  \paragraph{Central claim (renormalized IR locality).}
    In four dimensions, for a minimal Laplace-type operator and in the weak-curvature regime where all dimensionless
    curvature invariants satisfy \(R\,\ell_{\mathrm{sp}}^2\ll 1\), the infrared-dominant \emph{local}
    two-derivative part of the
    renormalized metric variation \(\delta S_\Pi^{\mathrm ren}/\delta g^{\mu\nu}\) is proportional to the Einstein tensor
    \(G_{\mu\nu}\).
    The associated coefficient defines an induced Newton coupling whose scaling is set by the spectral cutoff scale
    \(\ell_{\mathrm{sp}}\).
    The Einstein term arises from renormalization of the geometric sector, via the counterterm associated with the
    \(a_2\) heat-kernel pole, in the standard sense of induced gravity.
    The structural necessity of non-injectivity for any genuinely emergent description is established
    in~\cite{Beau2026l}.

    The present paper is a middle link in a derivation chain.
    The Laplace-type operator $A_g = -\nabla^2_g$ is taken here as the starting point;
    its derivation from the admissibility constraints of the Cosmochrony programme
    via the large-$q$ Mosco limit of the admissible fibre is addressed in~\cite{Beau2026Q5a}.
    The present paper establishes the infrared consequence: any such operator yields,
    after renormalization, a dominant Einstein response.

    Operationally, this result identifies a robust hierarchy.
    Local four-derivative contributions (quadratic curvature variations) and non-local form factors are suppressed in the
    infrared by additional powers of \(\ell_{\mathrm{sp}}^2\), while the conformal anomaly remains as a distinct
    ultraviolet-sensitive
    contribution whose presence is tracked explicitly.

    The remainder of the paper is organized as follows.
    In Section~\ref{sec:projective-entropy-covariant}, we define the covariant projective entropy functional and set up the
    minimal Laplace-type operator choice.
    In Section~\ref{sec:hk-sdw}, we summarize the heat-kernel expansion and the role of Seeley--DeWitt coefficients in the
    pole structure of zeta regularization.
    In Section~\ref{sec:metric-variation}, we define the renormalized spectral stress tensor and decompose it by derivative
    order into Einstein, quadratic-curvature, non-local, and anomaly pieces.
    In Section~\ref{subsec:ir-hierarchy}, we formulate the infrared hierarchy as a covariant derivative expansion and
    state the dominance criterion \(R\,\ell_{\mathrm{sp}}^2\ll 1\).
    In Section~\ref{sec:newton-identification}, we discuss the induced-gravity interpretation and the identification of the
    effective Newton coupling.
    We conclude in Section~\ref{sec:discussion} with limitations and outlook.
